\section{Visão Geral do Projeto}

Este trabalho examina a reescrita do Patch-Hub, uma ferramenta de código aberto desenvolvida sob a organização KWorkflow (kw), principalmente por estudantes da USP. O Patch-Hub é uma interface baseada em terminal para navegar e interagir com patches em lore.kernel.org, suportando o kernel Linux e projetos relacionados. Ele agiliza o fluxo de trabalho tradicional centrado em listas de discussão usado no desenvolvimento do kernel.

O Patch-Hub pertence ao projeto mais amplo KWorkflow, um Sistema de Fluxo de Trabalho de Automação de Desenvolvedor que reduz a sobrecarga de configuração do kernel e fornece ferramentas para tarefas do dia a dia. Reescrever o Patch-Hub usando o Modelo de Atores visa tanto melhorar a ferramenta quanto demonstrar os benefícios de arquiteturas baseadas em atores para aplicações centralizadas em Rust.

Este projeto é um caso de estudo adequado porque:

\begin{itemize}
    \item Tem um escopo e conjunto de funcionalidades claros
    \item Apresenta desafios arquiteturais bem adequados ao Modelo de Atores
    \item É mantido ativamente e usado pela comunidade do kernel
    \item Permite medição concreta do impacto arquitetural
\end{itemize}
